% The supplement's opening paragraph (adr/0007, adr/0008; RELEASE.md, "A paper with a supplement").
This supplement belongs to the paper of the same title. It holds the theory of pyramids on the line,
of which the paper prints in its Section~\ref{sec:line-theory} the part that the plane uses most.
Section~\ref{supp:strata} states the strata on the line in the line's own terms, with their proofs.
It works with the folded \Levy{} data and the operator axioms. At $d = 1$ its statements are
statements of Section~\ref{ch:pyramid} of the paper. Section~\ref{ch:line-appendix} treats the
discrete stratum on the line. It covers the gauge, the parametrization of one octave, the surface
measure and the knot-atom cascades, phase entropy, three clusters of examples, and the rays of the
classification on refining lattices.

Some proofs of the paper cite statements of this supplement. Those of
Section~\ref{sec:line-theory} read the exponent data on the line and discrete covariance there
(Lemmas~\ref{lem:line-exponent-data} and~\ref{lem:discrete-covariance-line}). Those of
Section~\ref{ch:matern} read the \Matern{} pyramid on the line: Lemma~\ref{lem:laplace-atoms-admissible},
Theorems~\ref{thm:matern-pyramid-line} and~\ref{thm:matern-pyramid-line-noncreation} and
Proposition~\ref{prop:ladder-non-creation-line}. The proof of Theorem~\ref{thm:matern-pyramid}, which
constructs the \Matern{} pyramid in $\RR^d$, rests on Theorem~\ref{thm:matern-pyramid-line}. A proof
of Section~\ref{ch:price} cites Proposition~\ref{prop:classical-kernels-no-exact-cascade}. The gauge
(Lemma~\ref{lem:gauge-group}) is cited in the paper's prose only. The
statements, sections and equations of the supplement are numbered S1, S.1, and so on. A number
without an S, such as Theorem~\ref{thm:lattice-strata}, is a number of the paper, and the notation
is that of the paper (Section~\ref{subsec:cited-notation}). Every statement here has a
machine-checked formal counterpart, under the trust base of Section~\ref{sec:trust} of the paper,
in the form Appendix~\ref{app:verification} describes.
